\begin{abstract}
    El abandono de clientes (\emph{churn}) afecta la continuidad de los ingresos en empresas de
    telecomunicaciones y motiva estrategias tempranas de retención. Este trabajo aborda su
    detección como una tarea de clasificación supervisada binaria sobre \emph{Telco Customer Churn},
    un conjunto público de ejemplo de IBM con 7043 clientes, 19 variables predictoras y una tasa de
    abandono del 26.54\%. Tras el análisis exploratorio y la revisión de calidad de datos, se
    implementó un protocolo reproducible con partición estratificada
    (60\% entrenamiento, 20\% validación y 20\% prueba), imputación de \var{TotalCharges} con cero y
    codificación \emph{one-hot}. Se comparó un clasificador de referencia (\emph{Dummy}) frente a
    una regresión logística (L2, $C=0.1$) y un árbol de decisión (\var{min\_samples\_leaf=20}),
    ambos con ponderación de clases (\var{balanced}). En el conjunto de prueba, la regresión
    logística alcanza \emph{Recall} de 0.781, F1-score de 0.617 y ROC-AUC de 0.841, superando al
    árbol en \emph{Recall} (0.722) y ROC-AUC (0.818), aunque sin diferencia concluyente en F1 según
    intervalos \emph{bootstrap} del 95\%. Ambos modelos identifican el contrato, la antigüedad
    y el tipo de internet como características predictivas relevantes, aunque comparten dificultades para
    clasificar abandonos en perfiles similares a los clientes que permanecen.
\end{abstract}

\begin{IEEEkeywords}
    Abandono de clientes, telecomunicaciones, clasificación binaria, regresión logística, árbol de
    decisión, desbalance de clases.
\end{IEEEkeywords}
